problem:  
Let \( X \) be a Calabi–Yau manifold admitting an (in general conjectural) special Lagrangian torus fibration \( \pi: X \to B \) with discriminant locus \( \Delta \subset B \), and let \( B_0 = B \setminus \Delta \) be the smooth locus endowed with an integral affine structure induced by \( X \). According to the Gross–Siebert program, one can associate to \( B_0 \) a *scattering diagram* \( \mathfrak{D}_A \) built from counts of Maslov index‑two holomorphic discs bounding Lagrangian fibers, encoding the A‑model instanton corrections. Similarly, one can construct another scattering diagram \( \mathfrak{D}_B \) from the B‑model side via degenerations of complex structures and log Gromov–Witten invariants of the mirror degeneration \( X^\vee_t \to \mathrm{Spec}\,\mathbb{C}[[t]] \).  

The two constructions \( \mathfrak{D}_A \) and \( \mathfrak{D}_B \) are each defined on integral affine manifolds with singularities, but their data a priori arise from distinct geometric input.  

**Problem:**  
Formulate and determine whether there exists a *canonical equivalence of scattering diagrams*  
\[
\Phi: (\mathfrak{D}_A, B_0) \longrightarrow (\mathfrak{D}_B, B_0^\vee)
\]  
up to gauge equivalence of wall structures, such that:  
1. The induced piecewise‑linear map on base affine manifolds identifies the underlying affine structures up to a Legendre transform;  
2. Under this identification, the wall‑crossing automorphisms determined by disc counts on the A‑side and log Gromov–Witten invariants on the B‑side coincide term‑by‑term after appropriate mirror coordinate change;  
3. The equivalence is *universal*, i.e., depends only on the underlying combinatorial data of the tropical manifold and not on the specific symplectic or complex realization.  

In essence, does mirror symmetry between \( X \) and \( X^\vee \) induce a canonical isomorphism between their *instanton‑corrected tropical scattering diagrams*, providing a purely tropical manifestation of mirror duality?  

Why is it a "good" problem:  
This formulation corrects vague aspects of the earlier version by grounding the question in well‑defined objects—scattering diagrams as constructed in the Gross–Siebert program—where equivalence up to gauge transformation and wall‑crossing automorphisms are rigorously defined. The problem now asks whether mirror symmetry admits a completely tropical description: a canonical identification between the A‑model and B‑model scattering diagrams encoding all instanton corrections. Such a correspondence, if it exists, would bridge enumerative symplectic geometry (disc counts) and logarithmic complex geometry (log Gromov–Witten invariants) through tropical wall structures alone.  
Conceptually, this problem probes whether the deep analytic content of mirror symmetry can be captured via combinatorial tropical data, achieving a synthesis of symplectic geometry, complex algebraic geometry, and tropical combinatorics. Its resolution would provide a foundation for a global, model‑independent theory of mirror symmetry and advance the unification of enumerative and tropical geometry—fulfilling all the criteria of a simple, profound, and cross‑disciplinary research problem.